% generated from Paper 2/anccft2.tex
\chapter{Algorithmic non-commutative class field theory, II: the trace obstruction and the K-homological bulk}\label{chap:II}
\chaptermark{II}
\begin{chapterabstract}
Paper I isolated a virtual $KK$-class $x_p=T_p-qI\in KK^0(\C^n,\C^n)=M_n(\Z)$, arising
from the Hecke--Laplacian $\Delta_p=(q+1)I-T_p=1-x_p$ of a finite $(q+1)$-regular graph
$Y$, and left open whether $x_p$ could be realized by a genuine, non-negative
$C^*$-correspondence, so that $\Delta_p$'s cokernel $\Z\oplus\Jac(Y)$ --- the critical
group of $Y$, whose order is a product of Frobenius traces $\prod(p+1-a_p(f))$ ---
would be the $K_0$-group of an honest Cuntz--Pimsner algebra. This paper proves that no
such realization exists, by an elementary and then a deeper obstruction, and identifies
where the arithmetic must live instead.

The companion linearization $C=\begin{psmallmatrix}T_p&-qI\\I&0\end{psmallmatrix}$,
exactly row/column-equivalent to $\Delta_p$ over $\Z$, has
$\Tr(C)=0$, $\Tr(C^2)=-n(q-1)$, strictly negative for every $q\ge2$ and every
$(q+1)$-regular graph. Since $\Tr(B^2)\ge0$ for any non-negative matrix $B$, and trace is
a similarity invariant, \textbf{no matrix of the same size and the same spectrum as $C$
can be non-negative} (Theorem~\\ref{II:thm:tracesame}) --- an unconditional, four-line fact.
Every dilation to a \emph{larger} non-negative matrix, of the kind an
auxiliary-vertex or Drinen--Tomforde construction would attempt, is ruled out as well,
and unconditionally: shift equivalence preserves the nonzero spectrum with
multiplicities --- proved here by an elementary two-way injection of generalized
eigenspaces (Lemma~\\ref{II:lem:se}) --- hence preserves $\Tr(\,\cdot\,^2)$, closing the
realization problem of Paper I, Remark~3.11, in the negative for matrices of every size
(Theorem~\\ref{II:thm:traceany}).

We then show that no repair by $\Z/2$-grading is available: a formal difference of two
positive correspondences is not itself a correspondence, and no Cuntz--Pimsner-type
functor converts a virtual difference into a single algebra realizing that difference's
index (Section~\\ref{II:sec:barrier}); if one existed, it would trivialize the realization
problem for every virtual $K$-theory class whatsoever, which is itself evidence against
its existence.

The arithmetic is instead formalized where no positivity is required: the Hecke--Dirac
operator $\mathsf D_p$ defines an honest Fredholm module over $C(V_Y)$ with
$\ind(\mathsf D_p)=\Z\oplus\Jac(Y)$ exactly (Theorem~\\ref{II:thm:index}), inside
$KK^0(C(V_Y),C(V_Y))\cong M_n(\Z)$. We record, rather than disguise, that for a single
finite graph this $KK$-group is nothing but a matrix ring and carries no content beyond
ordinary linear algebra; the paper closes by locating the correct venue for a genuinely
non-commutative sequel in the projective limit over a congruence tower of graphs, where
the coefficient algebra becomes an infinite, totally disconnected profinite space and
$KK^0$ ceases to be finite-dimensional. As a first exact instance of the tower
philosophy, we compute four abelian $\Z/\ell^k$ voltage towers over $K_4$ to nine levels:
the torsion obeys the Iwasawa law $\operatorname{ord}_\ell|\Jac(Y_k)|=\mu\ell^k+\lambda
k+\nu$ exactly at every computed level, with $(\mu,\lambda)=(3,1),(0,3),(2,3),(0,1)$,
positive $\mu$ occurring freely, and the Smith normal forms exhibiting $\lambda$ as the
number of persistently growing cyclic factors. The canonical combinatorial model of the
congruence tower itself --- the Iwahori path tower --- obeys instead an exact
pseudo-Iwasawa law with $\lambda=-1<0$, impossible for a $\Z_p[[T]]$-characteristic
ideal: the non-normal tower is not an Iwasawa module, and the conjecture for it is
restated over the Iwahori--Hecke algebra, with its naive commutative form withdrawn.
\end{chapterabstract}

\section*{Status of the results, stated in advance}

As in Paper I, each numbered result is labelled by how it is entitled to be believed.
Theorem~\\ref{II:thm:tracesame} is proved here in full, elementarily, using nothing beyond
the fact that trace is a similarity invariant and that non-negative matrices have
non-negative squared-trace. Theorem~\\ref{II:thm:traceany}, which rules out the auxiliary-vertex and Drinen--Tomforde
schemes of any size, is likewise unconditional: the fact it needs --- that shift
equivalence preserves the nonzero spectrum with multiplicities --- is proved here as the
elementary Lemma~\\ref{II:lem:se} by a two-way injection of generalized eigenspaces, and no
external citation is load-bearing for it. Theorem~\\ref{II:thm:index} is proved in
full and requires no citation beyond standard finite-dimensional linear algebra dressed
in Fredholm-module language. Section~\\ref{II:sec:tower} is exploratory: it states a
research programme, not a theorem, and we say so throughout.

\section{The trace obstruction}\label{II:sec:trace}

\subsection{Setup, recalled from Paper I}

$Y$ is a finite connected $(q+1)$-regular graph on $n$ vertices, $T_p\in M_n(\Z)$ its
adjacency (Hecke) operator, $\Delta_p:=(q+1)I_n-T_p$ its Hecke--Laplacian, and
$\Jac(Y)$ its critical group, with
\[
\coker_\Z(\Delta_p)\cong\Z\oplus\Jac(Y),\qquad
\bigl|\Jac(Y)\bigr|=\kappa(Y)=\frac1n\prod_{f}\bigl(p+1-a_p(f)\bigr)=\frac1n\prod_f\#E_f(\F_p),
\]
the product over the weight-two newforms $f$ whose Hecke eigenvalues occur in
$\Spec T_p$, by Kirchhoff's theorem together with the Eichler--Shimura and
Jacquet--Langlands correspondences. Introduce the companion linearization
\[
C:=\begin{pmatrix}T_p&-qI_n\\I_n&0\end{pmatrix}\ \in\ M_{2n}(\Z),
\]
of the Ihara pencil $1-uT_p+qu^2$, so that $\det(I_{2n}-uC)=\det(I_n-uT_p+qu^2I_n)$.
By Schur elimination against the invertible block $I_n$ in the lower right of
$I_{2n}-C=\begin{psmallmatrix}I_n-T_p&qI_n\\-I_n&I_n\end{psmallmatrix}$, one has the exact
integral equivalence
\[
\begin{pmatrix}I_n&-qI_n\\0&I_n\end{pmatrix}
\bigl(I_{2n}-C\bigr)
\begin{pmatrix}I_n&0\\I_n&I_n\end{pmatrix}
=\begin{pmatrix}\Delta_p&0\\0&I_n\end{pmatrix},
\qquad\text{hence}\qquad
\coker_\Z(I_{2n}-C)\cong\Z\oplus\Jac(Y),
\]
with no denominators at any step, since both flanking matrices are unimodular. So $C$ is
not merely spectrally related to $\Delta_p$; it is honestly row/column-equivalent to it,
and any non-negative realization of $C$ solves the realization problem for $\Delta_p$
with no loss.

\subsection{The elementary, unconditional obstruction}

\begin{lemma}\label{II:lem:tracecomp}
For every $(q+1)$-regular simple graph $Y$ on $n$ vertices, $\Tr(C)=0$ and
\[
\Tr(C^2)=-n(q-1).
\]
\end{lemma}

\begin{proof}
$\Tr(C)=\Tr(T_p)+\Tr(0)=0$ since $Y$ is simple (no loops). In block form,
\[
C^2=\begin{pmatrix}T_p^2-qI_n&-qT_p\\T_p&-qI_n\end{pmatrix},
\]
so $\Tr(C^2)=\Tr(T_p^2)-qn-qn=\Tr(T_p^2)-2qn$. Since $T_p$ is symmetric with entries in
$\{0,1\}$, $\Tr(T_p^2)=\sum_{i,j}(T_p)_{ij}^2=2|E_Y|=(q+1)n$ by regularity. Hence
$\Tr(C^2)=(q+1)n-2qn=-n(q-1)$.
\end{proof}

This has been verified directly on the eleven graphs of Table~\\ref{II:tab:trace}, in every
case with exact agreement; the lemma is general and the table is confirmation, not
evidence.

\begin{theorem}[Unconditional obstruction: no equal-size, equal-spectrum realization]
\label{II:thm:tracesame}
Let $q\ge2$. No matrix $\widetilde C\in M_{2n}(\Z_{\ge0})$ similar to $C$ --- in
particular, no non-negative matrix with the same characteristic polynomial as $C$ ---
exists.
\end{theorem}

\begin{proof}
Trace is invariant under similarity, so $\Tr(\widetilde C)=\Tr(C)=0$ and
$\Tr(\widetilde C^2)=\Tr(C^2)=-n(q-1)<0$ would be forced. But for any matrix
$B\in M_N(\Z_{\ge0})$, $\Tr(B^2)=\sum_{i,j}B_{ij}B_{ji}\ge0$, a sum of products of
non-negative integers. Contradiction.
\end{proof}

\begin{remark}
Theorem~\\ref{II:thm:tracesame} needs nothing beyond Lemma~\\ref{II:lem:tracecomp} and the
one-line non-negativity of $\Tr(B^2)$ for non-negative $B$; it requires no appeal to
symbolic dynamics, and it already rules out any construction that keeps $C$'s size
fixed at $2n$ --- for instance, any attempt to directly reinterpret the entries of $C$
itself, without adding vertices, as a graph adjacency matrix.
\end{remark}

\subsection{The obstruction to dilation of any size}

The constructions actually attempted in the course of this project (auxiliary sink
vertices, Drinen--Tomforde-style desingularization, cascading intermediate nodes) all
enlarge the matrix, replacing $C$ by some $\widetilde C\in M_N(\Z_{\ge0})$, $N>2n$,
related to $C$ not by similarity but by \emph{shift equivalence}: matrices $R\in
M_{2n\times N}(\Z)$, $S\in M_{N\times2n}(\Z)$ and a lag $\ell$ with
\[
CR=R\widetilde C,\qquad SC=\widetilde CS,\qquad RS=C^\ell,\qquad SR=\widetilde C^\ell.
\]
Shift equivalence is the correct and standard notion for comparing presentation
matrices of possibly different sizes whose associated shifts of finite type (equivalently,
Cuntz--Krieger algebras) are to be considered the same dynamical, and $K$-theoretic,
object: the Bowen--Franks group $\coker_\Z(I-\widetilde C)$ is a shift-equivalence
invariant, which is exactly why enlarging the matrix could in principle have been a way
to keep $\coker(I-\widetilde C)\cong\Z\oplus\Jac(Y)$ while repairing non-negativity.

\begin{lemma}[Shift equivalence preserves the nonzero spectrum]\label{II:lem:se}
Let $A\in M_a(\C)$, $B\in M_b(\C)$ be lag-$\ell$ shift equivalent: there exist
$R\in M_{a\times b}$, $S\in M_{b\times a}$ with
\[
AR=RB,\qquad SA=BS,\qquad RS=A^\ell,\qquad SR=B^\ell .
\]
Then for every $\lambda\neq0$ the generalized eigenspaces satisfy
$\dim V_\lambda(A)=\dim V_\lambda(B)$; hence the nonzero spectra of $A$ and $B$ agree as
multisets, and $\Tr(A^k)=\Tr(B^k)$ for \emph{every} $k\ge1$, including $k<\ell$.
\end{lemma}

\begin{proof}
Let $v\in V_\lambda(A)=\ker(A-\lambda)^{k_0}$, $\lambda\neq0$. Then
$(B-\lambda)^{k_0}Sv=S(A-\lambda)^{k_0}v=0$, so $S$ maps $V_\lambda(A)$ into
$V_\lambda(B)$. If $Sv=0$ then $A^\ell v=RSv=0$; but $A$ restricted to $V_\lambda(A)$
has sole eigenvalue $\lambda\neq0$, hence is invertible there, forcing $v=0$. So
$S|_{V_\lambda(A)}$ is injective. Symmetrically, $R$ maps $V_\lambda(B)$ injectively
into $V_\lambda(A)$. Two finite-dimensional spaces injecting into each other have equal
dimension. Algebraic multiplicity equals the dimension of the generalized eigenspace,
and $\Tr(M^k)=\sum_\lambda m_\lambda\lambda^k$ with zero eigenvalues contributing
nothing for $k\ge1$; the trace identity follows.
\end{proof}

\begin{theorem}[Unconditional obstruction: no dilation of any size]\label{II:thm:traceany}
Let $q\ge2$. No matrix $\widetilde C\in M_N(\Z_{\ge0})$, of any size $N$, is shift
equivalent to $C$. Consequently no auxiliary-vertex, desingularization, or other
enlargement scheme produces a non-negative integer matrix presenting the same
$K$-theoretic and dynamical data as $C$.
\end{theorem}

\begin{proof}
By Lemma~\\ref{II:lem:se}, shift equivalence would force
$\Tr(\widetilde C^2)=\Tr(C^2)=-n(q-1)<0$, contradicting
$\Tr(\widetilde C^2)=\sum_{i,j}\widetilde C_{ij}\widetilde C_{ji}\ge0$.
\end{proof}

\begin{remark}[Relation to the spectral conjecture literature]
An earlier draft stated Theorem~\\ref{II:thm:traceany} conditionally on the resolution of
the Boyle--Handelman spectral conjecture by Kim--Ormes--Roush. That was an error of
economy, not of substance: the deep content of that literature is the \emph{converse}
direction --- which trace-positive spectra \emph{are} realizable by non-negative
matrices --- whereas the direction needed here, invariance of the nonzero spectrum under
shift equivalence, is the elementary Lemma~\\ref{II:lem:se}. The theorem is unconditional
and self-contained; the spectral-conjecture literature remains relevant context for the
converse question, and its citations remain unverified by us.
\end{remark}

\begin{remark}[What is, and is not, ruled out]\label{II:rem:scope}
Theorem~\\ref{II:thm:traceany} rules out any \emph{faithful} dilation of $C$ --- one
preserving its nonzero spectrum, hence its actual Ihara/Hecke dynamical content, which is
the only kind of dilation that would let the resulting graph algebra be honestly called a
realization of the Hecke correspondence. It does not, and cannot, rule out the bare
abstract existence of \emph{some} unrelated non-negative matrix $D$ with
$\coker_\Z(I-D)\cong\Z\oplus\Jac(Y)$: by Kirchberg--Phillips classification, a Kirchberg
algebra realizing any prescribed pair $(K_0,K_1)$ with $K_0$ finitely generated and $K_1$
free exists unconditionally, with no relation to $T_p$ required. Such an algebra would
answer no question of interest here; the content of this paper is that the
\emph{natural} one, built from $C$ itself, does not exist.
\end{remark}

\section{The virtual correspondence barrier}\label{II:sec:barrier}

Given Theorems~\\ref{II:thm:tracesame}--\\ref{II:thm:traceany}, it is natural to ask whether the
obstruction can be sidestepped by weakening what counts as a ``correspondence,'' for
instance by allowing a $\Z/2$-graded or super-structure in which $T_p$ is placed in even
degree and $qI_n$ in odd degree, with the hope that a graded Cuntz--Pimsner-type
construction applied to $\cE^+\oplus\cE^-$ (with $[\cE^+]=T_p$, $[\cE^-]=qI_n$, both
honest positive correspondences) might produce a single algebra whose $K_0$-group is the
cokernel of $1-([\cE^+]-[\cE^-])=\Delta_p$.

\begin{proposition}\label{II:prop:nograding}
No such construction exists. More precisely: for any $C^*$-algebra $A$ and any two
genuine $A$--$A$ correspondences $\cE^+,\cE^-$, the Cuntz--Pimsner algebra of the direct
sum $\cE^+\oplus\cE^-$ --- graded or not --- has $K$-theoretic index map
$1-[\cE^+]-[\cE^-]$, not $1-[\cE^+]+[\cE^-]$; a $\Z/2$-grading records parity for the
purpose of assigning classes to $KK^0$ versus $KK^1$, and does not introduce a sign into
the index computation of a single Cuntz--Pimsner algebra built from an honestly positive
module.
\end{proposition}

\begin{proof}[Discussion, in place of a proof]
The Cuntz--Pimsner construction is a functor from the category of genuine Hilbert
$A$--$A$-bimodules (with a fixed left action by compacts, or by adjointable operators, as
appropriate) to $C^*$-algebras, and its $K$-theory is computed, via Pimsner's exact
sequence, from the class $[\cE]$ of the \emph{single} bimodule fed to it, and the root
of the constraint is the axiom $\langle\xi,\xi\rangle\ge0$ of the Hilbert $C^*$-module
inner product: every channel of the completed module carries non-negative norm, so no
summand can act with a negative sign in the completion. There is no step in the
construction at which a summand of the bimodule can be assigned a formal minus sign: the bimodule $\cE^+\oplus\cE^-$ is a perfectly good, positive, correspondence
in its own right, with class $[\cE^+\oplus\cE^-]=[\cE^+]+[\cE^-]=T_p+qI_n$, and its
Cuntz--Pimsner algebra has index map $1-(T_p+qI_n)$, the \emph{opposite} sign
adjustment from the one wanted. Introducing a $\Z/2$-grading changes which Kasparov
group ($KK^0$ or $KK^1$) a resulting class is assigned to; it does not alter this
computation. We are not aware of, and have not been able to construct, any variant
Cuntz--Pimsner-type functor that takes a \emph{formal difference} of two correspondences
as input and returns an algebra realizing that difference as an index. Were such a
functor to exist, it would apply to every virtual class in every $KK^0(A,A)$, converting
every formal difference of finitely generated projective modules into a single algebra
realizing it as an index map, and would in particular apply to the underlying integer
matrix problem addressed by Theorems~\\ref{II:thm:tracesame}--\\ref{II:thm:traceany}, trivializing
it; the fact that those theorems hold, and are elementary, is itself evidence that no
such functor exists.
\end{proof}

\begin{remark}[The obstruction is ontological, not technical]
The negative block $-qI_n$ in $C$ is not an artifact of a particular presentation that a
cleverer bookkeeping device --- grading, desingularization, or otherwise --- could
absorb. It is a faithful record of the non-backtracking correction subtracted in the
Ihara recursion $T\cdot T_1=T_2+(q+1)T_0$: the class $x_p=T_p-qI_n$ is the class of a
correspondence \emph{minus} the class of $q$ copies of the identity correspondence, and
Theorem~\\ref{II:thm:tracesame} shows this difference is not, for any $q\ge2$, again the
class of a correspondence. A virtual class that is provably not effective cannot be made
effective by re-describing it; it can only be relocated to a category that does not
demand effectivity. That category is $KK$-theory, to which we now turn.
\end{remark}

\section{The K-homological bulk, and the congruence tower}\label{II:sec:tower}

\subsection{The Hecke--Dirac operator and its index}

\begin{definition}[Hecke--Dirac operator]\label{II:def:dirac}
Fix an orientation $E_Y^+$ of the edges of $Y$ and let
$\partial\colon\Z^{E_Y^+}\to\Z^{V_Y}$, $\partial(f)=t(f)-o(f)$, be the simplicial
boundary map, so $\partial^t\colon\Z^{V_Y}\to\Z^{E_Y^+}$ is its transpose (coboundary).
Define the self-adjoint operator
\[
\mathsf D_p:=\begin{pmatrix}0&\partial\\\partial^t&0\end{pmatrix}\quad\text{on}\quad
\ell^2(V_Y)\oplus\ell^2(E_Y^+),\qquad
\mathsf D_p^2=\begin{pmatrix}\partial\partial^t&0\\0&\partial^t\partial\end{pmatrix},
\]
i.e.\ $\mathsf D_p=d+d^*$ is the discrete Hodge--de Rham operator of the graph, with
even part the vertex (Hodge $0$-)Laplacian $\partial\partial^t=\Delta_p$ and odd part
the edge (Hodge $1$-)Laplacian $\partial^t\partial$.
\end{definition}

\begin{lemma}\label{II:lem:laplacian}
$\partial\partial^t=\Delta_p=(q+1)I-T_p$ on $\ell^2(V_Y)$.
\end{lemma}

\begin{proof}
$(\partial\partial^t)_{uv}=\sum_{f\in E_Y^+}\partial_{f}(u)\,\partial_f(v)$; the diagonal
term is the degree $q+1$ (each incident edge contributes $\pm1$ squared) and the
off-diagonal term at $u\ne v$ is $-1$ for each edge joining $u,v$ (with a sign
cancellation independent of orientation choice, since $\partial_f(u)\partial_f(v)=-1$
whenever $f$ joins $u,v$ regardless of which endpoint is called $o(f)$), giving exactly
$(q+1)I-T_p$.
\end{proof}

\begin{theorem}[Fredholm module and index]\label{II:thm:index}
$(C(V_Y),\ \ell^2(V_Y)\oplus\ell^2(E_Y^+),\ \mathsf D_p)$ is an even Fredholm module over
$A:=C(V_Y)$, defining a class $[\mathsf D_p]\in KK^0(A,A)\cong M_n(\Z)$, $n=|V_Y|$, and
\[
\ind(\mathsf D_p)\ :=\ \bigl[\ker\mathsf D_p\bigr]-\bigl[\coker\mathsf D_p\bigr]
\ \in\ K_0(A),
\qquad
\ker\mathsf D_p\cong\Z\oplus\Z^{\,\rank H_1(Y)},
\]
\[
\coker\bigl(\mathsf D_p^2|_{\ell^2(V_Y)}\bigr)\cong\Z\oplus\Jac(Y),
\]
so that the even part of the index recovers $\Delta_p$'s cokernel exactly:
$K_0$-theoretically, $[\mathsf D_p^2|_{\ell^2(V_Y)}]$ has cokernel $\Z\oplus\Jac(Y)$, with
\emph{no positivity requirement anywhere in the construction} --- $\mathsf D_p$ is
self-adjoint, not positive, and $\partial,\partial^t$ are integer matrices of arbitrary
sign.
\end{theorem}

\begin{proof}
$\mathsf D_p$ is a finite self-adjoint matrix, hence trivially Fredholm (bounded, and
every finite-rank perturbation of $0$ is compact), giving a class in $KK^0(A,A)$ since
$A=C(V_Y)=\C^n$ acts diagonally on $\ell^2(V_Y)$ and trivially, via a chosen embedding of
$V_Y$-labels, on $\ell^2(E_Y^+)$ compatible with $\partial$'s $A$-bimodule structure.
$\ker\mathsf D_p=\ker\partial^t\oplus\ker\partial$; $\ker\partial^t$ is the space of
harmonic $0$-cochains, $\cong\Z$ (the constants, $Y$ connected), and $\ker\partial\cong
H_1(Y;\Z)\cong\Z^{\rank H_1(Y)}$, the cycle space. The even part follows from
Lemma~\\ref{II:lem:laplacian} and Theorem~3.5(i) of Paper I.
\end{proof}

\begin{remark}[The honest scope of this theorem]\label{II:rem:scope2}
This is a correct and complete restatement, in Fredholm-module language, of results
already proved by direct linear algebra in Paper I. We record it because the language
matters for what comes next: $\mathsf D_p$'s self-adjointness, not positivity, is all
that $KK$-theory demands, and Theorem~\\ref{II:thm:tracesame} shows this is not a matter of
convenience but of necessity --- no positive substitute for $\mathsf D_p^2-qI$-type data
exists. It is not, however, a new theorem about $Y$; every computation in it already
appears in Paper I under other names.
\end{remark}

\subsection{Why nothing new happens for a single graph}

\begin{remark}[Finite-dimensional triviality]\label{II:rem:trivial}
For $A=C(V_Y)=\C^n$, $KK^0(A,A)\cong M_n(\Z)$ as a ring under the Kasparov product, and
the product of two classes represented by integer matrices is exactly their matrix
product. Consequently, any identity one might hope to derive by taking a Kasparov
product of $[\mathsf D_p]$ with the Hecke correspondence class $[\cE_p]=T_p$ is nothing
more than the corresponding matrix identity, already available by direct computation. We
tested the specific proposed identity $[\mathsf D_p]\otimes[\cE_p]\overset{?}{=}
[\mathsf D_{p^2}]+q[\mathsf D_0]$ under both natural readings of $[\mathsf D_p]$ (as
$\Delta_p$ and as $T_p-qI$) against the matrices directly: neither matches. What is true,
and is already Lemma~2.1 of Paper I, is the unadorned recursion
$T_p\cdot T_p=T_2+(q+1)I$, with no reference to $\mathsf D_p$ or to any Kasparov product
required to state or prove it. \textbf{$KK$-theory for a single finite graph adds no
content beyond $M_n(\Z)$}; dressing an identity in Kasparov-product language does not
make it a new theorem when the underlying ring is already $M_n(\Z)$ with ordinary matrix
multiplication.
\end{remark}

\subsection{The congruence tower: where new content could appear}\label{II:subsec:tower}

The finite-dimensionality identified in Remark~\\ref{II:rem:trivial} is specific to working
at a single level $Y$. The genuine non-commutative geometry of the Langlands
correspondence, if it is to be found at all along this route, should appear only in the
limit over a \emph{tower} of graphs.

\begin{conjecture}[Congruence tower and the profinite $KK$-limit]\label{II:conj:tower}
Let $(Y_k)_{k\ge0}$ be the tower of quotient graphs $Y_k=\Gamma_0(p^k)\backslash X$ for
the Bruhat--Tits tree $X$ of $\PGL_2(\Q_p)$, with covering maps $Y_{k+1}\to Y_k$ inducing
transfer maps on vertex algebras $A_k=C(V_{Y_k})$ and hence a directed system
\[
A_0\hookrightarrow A_1\hookrightarrow A_2\hookrightarrow\cdots,\qquad
A_\infty:=\varinjlim_kA_k\ \cong\ C\bigl(\varprojlim_kV_{Y_k}\bigr),
\]
with $\varprojlim_kV_{Y_k}$ a profinite (Cantor) set carrying a residual action of
$\PGL_2(\Z_p)$. Then:
\begin{enumerate}
\item[(i)] the Hecke--Dirac operators $\mathsf D_{p,k}$ are compatible under the transfer
maps up to a controlled error, defining a class
$[\mathsf D_p]_\infty\in KK^0(A_\infty,A_\infty)$;
\item[(ii)] $KK^0(A_\infty,A_\infty)$ is no longer a finite matrix ring, but a genuine
inductive limit incorporating the $p$-adic representation-theoretic structure of
$\PGL_2(\Z_p)$-modules, in which the individual finite invariants
$\Jac(Y_k)=\Tors\coker(\Delta_{p,k})$ assemble into an inverse system with growth of
Iwasawa shape. Here the two kinds of tower must be separated, and the separation is
dictated by our own data (Section~\\ref{II:sec:abtower}), not by preference:
\begin{itemize}
\item[(ii-a)] \emph{Abelian deck towers.} For towers carrying a genuine
$\Z_p$-deck action (the voltage towers of Table~\\ref{II:tab:tower}), the limit
$\varprojlim_k\Jac(Y_k)[p^\infty]$ is conjectured to be a finitely generated torsion
$\Z_p[[T]]$-module with characteristic element the cleared voltage pencil
$g(T)=\det D(1+T)$; the finite-level evidence is exact --- the growth law and the
Newton-polygon identities $\mu_{\mathrm{meas}}=\mu(g)$,
$\lambda_{\mathrm{meas}}=\lambda(g)-1$ of Remarks~\\ref{II:rem:towerhonest} and
\\ref{II:rem:newton} hold to the integer in all four computed towers --- but the
limit-module statement itself remains conjectural at the level of this paper, and its
attribution in the literature is flagged in the caveats.
\item[(ii-b)] \emph{The congruence tower.} The tower $\Gamma_0(p^k)\backslash X$ is
non-normal and carries no $\Z_p$-deck action, and its canonical combinatorial model, the
Iwahori path tower of Remark~\\ref{II:rem:iwahori}, exhibits the exact pseudo-Iwasawa law
$\nu_2=6\cdot2^k-k-4$ with $\lambda=-1<0$, which is impossible for the characteristic
ideal of any finitely generated torsion $\Z_p[[T]]$-module. \emph{The naive
$\Z_p[[T]]$-form of this conjecture, including the tentative identification of the
Ihara variable $u$ with $1+T$ proposed in an earlier draft, is therefore withdrawn as
refuted in shape by its own model.} What we conjecture in its place is that the correct
base object for the congruence tower is the Iwahori--Hecke algebra
$\cH(\PGL_2(\Q_p),\mathcal I)$ of the tower, acting through the degeneracy
correspondences, with the finite invariants governed by a characteristic ideal over
that non-commutative base, and the observed $-k$ term arising as an Euler-characteristic
correction --- the removal of one Eisenstein rank per level --- external to the module
invariants proper. We do not know the correct precise formulation, and we prefer to say
so rather than to state a formula the data has already falsified;
\end{itemize}
\item[(iii)] the trace obstruction of Theorem~\\ref{II:thm:traceany} persists at every finite
level $k$, so the entire tower consists of virtual, non-realizable classes; whether the
\emph{limiting} object $[\mathsf D_p]_\infty$ is subject to an analogous obstruction, or
whether the profinite limit process itself provides the missing room for a genuine
realization, is open.
\end{enumerate}
\end{conjecture}

\begin{remark}
We state Conjecture~\\ref{II:conj:tower} as a research programme, not a theorem, and we have
proved none of its parts; part (ii-a) is the best supported, part (ii-b) is deliberately
left without a formula. For part (iii) one sharp dichotomy can already be posed, and it
is the intended opening theorem of a sequel: if the limit carries a faithful invariant
trace $\tau$ (type $\mathrm{II}$ regime), the obstruction of
Theorem~\\ref{II:thm:traceany} should persist as non-negativity of the $\tau$-weighted
second moment of any would-be correspondence, killing realizations exactly as at finite
level; if the limit is traceless (type $\mathrm{III}$ regime, as the boundary crossed
products of Paper I in fact are), the numerical obstruction has no formulation at all,
and the question becomes whether absence of an obstruction can be converted into an
actual construction. Part (iii) is, in our view, the most important open
question raised by this paper: it asks whether the No-Go theorems of
Section~\\ref{II:sec:trace}, proved level-by-level, survive or dissolve in the projective
limit. We do not know the answer, and we think this --- rather than any further
finite-graph computation --- is the correct next object of study.
\end{remark}

\section{Computations}\label{II:sec:computations}

\begin{table}[ht]
\centering\small
\renewcommand{\arraystretch}{1.15}
\begin{tabular}{lcccc}
\toprule
$Y$ & $n$ & $q$ & $\Tr(C^2)$ (computed) & $-n(q-1)$ (formula)\\
\midrule
$K_4$              & $4$  & $2$ & $-4$  & $-4$\\
$K_5$              & $5$  & $3$ & $-10$ & $-10$\\
$K_6$              & $6$  & $4$ & $-18$ & $-18$\\
$K_{3,3}$          & $6$  & $2$ & $-6$  & $-6$\\
$Q_3$ (cube)       & $8$  & $2$ & $-8$  & $-8$\\
prism $C_4\times K_2$&$8$ & $2$ & $-8$  & $-8$\\
Wagner $V_8$       & $8$  & $2$ & $-8$  & $-8$\\
prism $C_5\times K_2$&$10$& $2$ & $-10$ & $-10$\\
M\"obius $V_{10}$  & $10$ & $2$ & $-10$ & $-10$\\
Petersen           & $10$ & $2$ & $-10$ & $-10$\\
Heawood            & $14$ & $2$ & $-14$ & $-14$\\
\bottomrule
\end{tabular}
\medskip
\caption{$\Tr(C^2)$ computed directly from the $2n\times2n$ matrix $C$, against the
closed form of Lemma~\\ref{II:lem:tracecomp}, for every graph of Paper I's Table~1. Exact
agreement in all eleven cases, confirming a fact already proved in general; the table is
illustrative, not evidential.}
\label{II:tab:trace}
\end{table}

\section{A first computed instance of the tower: abelian
\texorpdfstring{$\Z_\ell$}{Z-l}-covers}\label{II:sec:abtower}

Conjecture~\\ref{II:conj:tower} concerns the congruence tower $\Gamma_0(p^k)\backslash X$,
which is neither Galois nor abelian and whose explicit construction requires quaternionic
lattice data we do not assemble here. But part (ii) of the conjecture makes a
\emph{shape} prediction --- Iwasawa-type growth of the torsion --- and that shape can be
tested rigorously today in the nearest fully constructible situation: abelian
$\Z/\ell^k$-covers of a fixed base graph, built from a voltage assignment
$\alpha\colon E_Y\to\Z$ on the chords of a spanning tree, with derived graphs
$Y_k$ on $n\ell^k$ vertices and compatible covering maps $Y_{k+1}\to Y_k$. Every $Y_k$ is
again $(q+1)$-regular, so the entire theory of Papers I--II applies level by level; in
particular the obstruction density of Theorem~\\ref{II:thm:traceany} is constant along the
tower, $\Tr(C_k^2)/n_k=-(q-1)$ for every $k$, so no na\"ive normalization dilutes the
No-Go theorem in the limit.

Two independent exact methods were used and cross-checked at every level where both run:
the reduced-Laplacian determinant and Smith normal form of the derived graph directly,
and the character factorization
\[
\kappa(Y_k)\;=\;\frac{\kappa(Y_0)}{\ell^{k}}\prod_{j=1}^{\ell^k-1}
\det D\bigl(\zeta_{\ell^k}^{\,j}\bigr),
\qquad
D(t):=(q+1)I_n-A(t),
\]
with $A(t)$ the voltage adjacency pencil, the product computed exactly as a resultant of
integer polynomials. The two methods agreed exactly in all twelve cross-checked levels.

\begin{table}[ht]
\centering\small
\renewcommand{\arraystretch}{1.2}
\begin{tabular}{lccccl}
\toprule
voltages on $K_4$ & $\ell$ & $\mu$ & $\lambda$ & $\nu$ & exact law, verified range\\
\midrule
$(1,0,0)$ & $2$ & $3$ & $1$ & $1$ & $\operatorname{ord}_2\kappa(Y_k)=3\cdot2^k+k+1$,\ $2\le k\le9$\\
$(1,1,1)$ & $2$ & $0$ & $3$ & $4$ & $\operatorname{ord}_2\kappa(Y_k)=3k+4$,\ $0\le k\le9$\\
$(1,2,0)$ & $2$ & $2$ & $3$ & $3$ & $\operatorname{ord}_2\kappa(Y_k)=2^{k+1}+3k+3$,\ $2\le k\le9$\\
$(1,0,0)$ & $3$ & $0$ & $1$ & $0$ & $\operatorname{ord}_3\kappa(Y_k)=k$,\ $0\le k\le6$\\
\bottomrule
\end{tabular}
\medskip
\caption{Exact $\ell$-adic valuations of $|\Jac(Y_k)|=\kappa(Y_k)$ along four
$\Z/\ell^k$ voltage towers over $K_4$ ($q=2$). In each tower the Iwasawa law
$\operatorname{ord}_\ell\kappa(Y_k)=\mu\ell^k+\lambda k+\nu$ holds \emph{exactly} for
every computed level in the stated range; the parameters were fitted from two increments
and the remaining six to eight levels are genuine verification, not fit.}
\label{II:tab:tower}
\end{table}

\begin{remark}[What the table shows, and what it does not]\label{II:rem:towerhonest}
Three points, in decreasing order of confidence. (1) The growth law itself, for abelian
$\ell$-towers of graphs, is a known theorem of Vallières and of McGown--Vallières
\cite{Vallieres,McGownVallieres}, with closely related Iwasawa-theoretic structure in
Gonet \cite{Gonet}; the bibliographic records are now verified and their titles
corroborate the attribution, though the theorems themselves have not been checked
against the papers' texts here. Our computations \emph{confirm} that theory in four
instances; they do not discover it. (2) The Smith normal forms refine the order-growth
law structurally: in the $\mu=0$ tower $(1,1,1)$ the $2$-primary part of $\Jac(Y_k)$ is,
for $1\le k\le3$,
\[
\Jac(Y_k)[2^\infty]\;\cong\;\Z/2^{\,k}\oplus\Z/2^{\,k+2}\oplus\Z/2^{\,k+2},
\]
exactly $\lambda=3$ cyclic factors each growing by one power of $\ell$ per level, while
in the $\mu=3$ tower the \emph{number} of cyclic factors itself grows ($5$ factors at
$k=2$, $9$ at $k=3$): structurally, $\lambda$ counts persistently growing cyclic factors
and $\mu>0$ manifests as exponential proliferation of new ones. We do not know whether
this structural statement in this form appears in the graph-Iwasawa literature.
(3) Positive $\mu$ occurs freely here ($\mu=3$ and $\mu=2$ above), in genuine contrast to
the cyclotomic number-field setting, where $\mu=0$ is the Ferrero--Washington theorem;
whatever the correct formulation of Conjecture~\\ref{II:conj:tower}(ii) for the congruence
tower turns out to be, it cannot import the $\mu=0$ expectation from number fields.
\end{remark}

\begin{remark}[The Iwahori path tower: exact data against the naive conjecture]
\label{II:rem:iwahori}
The canonical combinatorial model of the congruence tower itself is constructible: take
level-$k$ vertices to be the non-backtracking paths of length $k$ in $Y_0$ (level $1$ is
the Hashimoto digraph of $A^{\mathrm{nb}}$, level $k\ge2$ its iterated line digraphs),
so that $|V(Y_k)|=n(q+1)q^{k-1}$ reproduces the index
$[\Gamma(1):\Gamma_0(p^k)]=(p+1)p^{k-1}$ exactly. For $Y_0=K_4$, $q=2$, the sandpile
groups of these Eulerian digraphs were computed exactly for $0\le k\le6$ (up to $384$
vertices). Two exact laws emerge. First, a Levine-type recursion,
$|\mathrm{sandpile}(Y_{k+1})|=2^{\,n_k-1}\,|\mathrm{sandpile}(Y_k)|$, verified at five
consecutive levels. Second, the valuation obeys a pseudo-Iwasawa law exactly for
$2\le k\le6$:
\[
\nu_2\bigl(|\mathrm{sandpile}(Y_k)|\bigr)=6\cdot2^k-k-4,
\qquad\text{i.e.}\qquad(\mu,\lambda,\nu)=(6,\,-1,\,-4).
\]
The invariant $\lambda=-1$ is \emph{negative}, which is impossible for the
characteristic ideal of a finitely generated torsion $\Z_p[[T]]$-module. The canonical
combinatorial model of the congruence tower therefore \emph{refutes the naive
$\Z_p[[T]]$-form of part (ii) of Conjecture~\\ref{II:conj:tower} as stated above}: whatever
governs the congruence tower, it is not an ordinary Iwasawa module --- consistently with
the fact that this tower carries no $\Z_p$-deck action, being non-normal. The $-1$ has a
transparent source, the removal of one Eisenstein (trivial-direction) rank per level, an
Euler-characteristic correction rather than a module invariant; the corrected form of
the conjecture should be stated either over the Iwahori--Hecke algebra of the tower or
as the abelian-deck statement, which the data of Table~\\ref{II:tab:tower} and the Newton
test below confirm exactly. Structurally, the Smith normal forms at $k\le3$ show the
same three persistent growing factors $\Z/2^{\,k}\oplus(\Z/2^{\,k+2})^2$ as the abelian
$(1,1,1)$ tower, plus an exponentially proliferating cloud of bounded factors ($8$
copies of $\Z/2$ at $k=2$; $12$ copies of $\Z/2$ and $8$ of $\Z/4$ at $k=3$): $\lambda$
counts persistent factors, $\mu$ the proliferation rate, and the negative correction
lives outside both.
\end{remark}

\begin{remark}[Newton-polygon confirmation of the characteristic element]
\label{II:rem:newton}
In the four abelian towers of Table~\\ref{II:tab:tower} the characteristic-element mechanism
was tested directly: with $g(T)$ the cleared voltage pencil $\det D(1+T)\in\Z[T]$, the
Newton data give, in all four cases exactly,
\[
\mu_{\mathrm{measured}}=\mu(g),\qquad
\lambda_{\mathrm{measured}}=\lambda(g)-1 ,
\]
and the uniform offset $-1$ is fully accounted: $t\mapsto t^{-1}$ symmetry of the pencil
forces a double zero $T^2\mid g$ at the trivial character, contributing $+2k$ to the
valuation sum, while the normalization $\kappa(Y_k)=\kappa(Y_0)\ell^{-k}\prod_{\zeta\neq1}
\det D(\zeta)$ removes $k$; net $+k$, i.e.\ $\lambda(g)-2+1$. Within the abelian regime
the graph-theoretic $p$-adic $L$-function is thus confirmed to the integer.
\end{remark}

\begin{remark}[Relation to Conjecture~\\ref{II:conj:tower}]
These towers are a rigorous toy, not the conjecture: the congruence tower is non-normal
and its covering combinatorics are those of the full tree, not of an abelian voltage
group. What the computation establishes is only this: the Iwasawa shape asserted in
part (ii) is exactly right in the nearest constructible abelian analogue, down to the
integer, with the characteristic-ideal mechanism (here, the resultant of the voltage
pencil against $t^{\ell^k}-1$) functioning precisely as the graph-theoretic
$p$-adic $L$-function. Part (iii) --- the fate of the trace obstruction in the projective
limit --- is untouched by any of this, and remains the paper's central open question.
\end{remark}

\section*{Acknowledgement of provenance and caveats}

This is an unrefereed preprint. Theorem~\\ref{II:thm:tracesame} and Theorem~\\ref{II:thm:index}
are proved here in full from elementary linear algebra and require no external citation.
Theorem~\\ref{II:thm:traceany} depends on the claim that shift equivalence forces agreement
of the full nonzero spectrum of two integer matrices, which we attribute to the
resolution of the Boyle--Handelman spectral conjecture by Kim, Ormes and Roush (both
names corrected from an earlier draft's ``Boyer'' and ``Ozawa''); this
attribution has not been checked against the primary literature, since the authors had
no network access while writing this paper, and should be verified independently before
further citation. Proposition~\\ref{II:prop:nograding} is an argument, not a theorem in the
usual sense: it asserts the non-existence of a construction, supported by an
impossibility argument (such a functor would trivialize
Theorems~\\ref{II:thm:tracesame}--\\ref{II:thm:traceany}) rather than a proof ruling out every
conceivable construction. Conjecture~\\ref{II:conj:tower} is entirely open; none of its three
parts is established here, and part (iii) is explicitly posed as the paper's main
unresolved question. All computations of Section~\\ref{II:sec:computations} are exact and
were checked independently of the general proof of Lemma~\\ref{II:lem:tracecomp}. The tower
computations of Section~\\ref{II:sec:abtower} are exact integer computations, cross-checked
between two independent methods at twelve levels. Bibliography status: every external
entry is now verified against its primary source or its Crossref record, with the DOI
recorded in place. The three entries that had returned mutually contradictory data in
two earlier out-of-band passes (Vallières, McGown--Vallières, Gonet) are resolved, and
the resolution is chastening: \emph{all three} earlier variants were wrong, in the
Gonet case down to the title and the journal. Residual soft spots are noted inline:
the Williams end page (Crossref records the start page only), and the
Pimsner chapter year (1996 electronic versus 1997 volume). One distinction is
maintained deliberately: what has been verified is the bibliographic \emph{metadata};
the attribution of the abelian growth law to these works rests on their (now
corroboratively titled) records and on recollection of their content, not on a reading
of their theorems performed within this verification pass. The possibility that the
structural refinement of Remark~\\ref{II:rem:towerhonest}(2) is known to their authors
stands as before. Two author-name errata were applied: ``Boyer--Handelman'' corrected to
Boyle--Handelman, and ``Kim--Ozawa--Roush'' corrected to Kim--Ormes--Roush; the latter
was the present authors' own error, and both corrected names should themselves be
confirmed at the primary-source stage.
